\chapter{The Toolbox: Graphs, Grids, Time, Motion and Uncertainty}
\label{ch:ch02}
\chaptermark{The Toolbox}  % short running head; the full title is too long next to the section mark
% Function names for pseudocode are prefixed with "Toolbox" to stay unique.
\SetKwFunction{ToolboxLazyPush}{LazyPush}
\SetKwFunction{ToolboxLazyPop}{LazyPop}

\begin{objectives}
\begin{itemize}
  \item Define graphs, grid graphs and 3D lattices precisely, including the cost of a diagonal move and the corner-cutting rule.
  \item Explain why a disk-shaped drone can be planned as a point once the obstacles have been inflated by a Minkowski sum.
  \item Tell a path, a plan and a trajectory apart; build the time-expanded graph with wait actions; compute path length, travel time, makespan, sum of costs and minimum separation.
  \item Write down the single- and double-integrator models of a drone with velocity and acceleration limits, and compute its stopping distance.
  \item Use the geometric primitives of later chapters: distance to a segment, ray--circle intersection, tangents to a circle and the time of closest approach.
  \item Read a covariance matrix as an ellipse, state Bayes' rule, and use the lazy-deletion priority queue and the hashed states on which every search in this book relies.
\end{itemize}
\end{objectives}

% ---------------------------------------------------------------------
\section{Why this matters for a drone swarm}
\label{sec:ch02-motivation}

Picture eight quadrotors inspecting the racks of a warehouse. A global planner has given each of them a route. Halfway through the mission a forklift-mounted drone that nobody planned for crosses the main aisle. To handle this scene the four layers of \cref{ch:ch24} need a shared vocabulary. The global planner needs a \emph{map} it can search: a graph, usually a grid. It also needs to know how \emph{big} a drone is, so that a route through a narrow gap is not a route into a shelf. It needs \emph{time}, because two drones may use the same aisle if they do so at different moments, and \emph{cost measures} to compare routes and whole multi-agent plans.

The local safety layer needs the drone's \emph{motion limits} (how fast, how hard it can brake) and a small kit of \emph{geometry}: when do two moving disks touch, and which velocities steer clear of an obstacle? The prediction layer describes the intruder with a mean and a covariance, so the planner must understand \emph{uncertainty}. Finally, everything has to run in a few milliseconds, which is a question of \emph{computation}: complexity, priority queues and hashing.

This chapter collects these tools in one place. It is the only chapter with no planning algorithm of its own---its single piece of pseudocode is the priority queue that all the others run on---and it fixes instead the definitions and conventions that every later chapter cites by number. Read it once, then use it as a reference; the drone box of \cref{sec:ch02-drone} says which layer uses what. Every formal object gets a numbered definition, and every number quoted in the text comes from one of the companion files \code{code/ch02\_toolbox.py}, \code{code/figures/gen\_ch02\_gaussian.py} and \code{code/figures/gen\_ch02\_double\_integrator.py}.


% ---------------------------------------------------------------------
\section{Graphs, grids and lattices}
\label{sec:ch02-graphs}

\begin{keyidea}
Every planner in this book searches a graph. The graph may be explicit (a road map stored in memory), implicit (a grid whose neighbors are computed on demand) or time-expanded (one copy of the map per time step). The search algorithms of Parts~II and~III do not care which; they only ask a vertex for its successors and their costs.
\end{keyidea}

\subsection{Graphs, paths and shortest paths}
\label{sec:ch02-graph-defs}

\begin{definition}[Graph, weighted graph]
\label{def:ch02-graph}
A \textbf{directed graph}\index{graph!directed} $G=(V,E)$ consists of a finite set $V$ of \textbf{vertices} (also called \emph{nodes}) and a set $E\subseteq V\times V$ of \textbf{edges}; the edge $(u,v)$ leads from $u$ to $v$. The graph is \textbf{undirected}\index{graph!undirected} if $(u,v)\in E$ implies $(v,u)\in E$; we then write $\set{u,v}$ for the pair of opposite edges and count it once. A \textbf{weighted graph}\index{graph!weighted} carries a cost function $c\colon E\to[0,\infty)$; $c(u,v)$ is the \textbf{edge cost}. In an undirected weighted graph we also require $c(u,v)=c(v,u)$, so that the pair $\set{u,v}$ carries a single cost. The \textbf{successors} of $u$ are $\Succ(u)=\set{v : (u,v)\in E}$ and its \textbf{predecessors} are $\Pred(u)=\set{v : (v,u)\in E}$; $\abs{\Succ(u)}$ is the out-degree of $u$, and the largest out-degree in the graph is the \textbf{branching factor} $b$.
\end{definition}

Edge costs in this book are non-negative, and almost always positive: a move on a grid costs its length, a wait costs one time step. Dijkstra's algorithm (\cref{ch:ch03}) needs exactly this non-negativity.

A graph can be stored in three ways. An \emph{adjacency matrix} keeps a $\abs{V}\times\abs{V}$ table of costs; it answers ``is there an edge from $u$ to $v$?'' in constant time but needs $\bigO{\abs{V}^2}$ memory. An \textbf{adjacency list}\index{adjacency list} keeps, for every vertex, the list of its successors with their costs; in \python this is a dictionary such as \code{\{'a': [('b', 1.0), ('c', 2.5)], 'b': [('c', 1.0)], 'c': []\}}, and it needs $\bigO{\abs{V}+\abs{E}}$ memory. An \emph{implicit graph}\index{graph!implicit} stores nothing at all: a successor function computes $\Succ(u)$ on demand. Grids are implicit graphs, and so are the time-expanded graphs of \cref{sec:ch02-time}. All the search code of this book is written against the implicit interface, a function \code{neighbors(v)} that returns pairs \code{(v', cost)}, because an adjacency list can always be wrapped in such a function.

\begin{definition}[Path, path cost]
\label{def:ch02-path}
A \textbf{path}\index{path} from $s$ to $\gamma$ in $G$ is a sequence $\pi=(v_0,v_1,\dots,v_k)$ of vertices with $v_0=s$, $v_k=\gamma$ and $(v_{i-1},v_i)\in E$ for $i=1,\dots,k$. Its \textbf{cost}\index{path!cost} is the sum of its edge costs,
\begin{equation}
\cost(\pi)=\sum_{i=1}^{k} c(v_{i-1},v_i),
\label{eq:ch02-path-cost}
\end{equation}
and $k$ is its number of edges. A path is \emph{simple} if no vertex appears twice. The one-vertex sequence $(s)$ is a path from $s$ to $s$ of cost $0$.
\end{definition}

\begin{definition}[Shortest path, distance]
\label{def:ch02-shortest-path}
A \textbf{shortest path}\index{shortest path} from $s$ to $\gamma$ is a path of minimum cost among all paths from $s$ to $\gamma$. Its cost is the \textbf{distance} $\dist(s,\gamma)=\min_\pi \cost(\pi)$; we set $\dist(s,\gamma)=\infty$ when no path exists, and $\dist(s,s)=0$. The distance from $s$ to $\gamma$ is also written $\hcost^*(s)$ when $\gamma$ is fixed (\cref{ch:ch04}).
\end{definition}

\begin{proposition}[Triangle inequality]
\label[proposition]{thm:ch02-triangle}
For all vertices $u,v,w$ of a graph with non-negative costs, $\dist(u,w)\le \dist(u,v)+\dist(v,w)$.
\end{proposition}
\begin{proof}
If either distance on the right is infinite there is nothing to show. Otherwise take a shortest path from $u$ to $v$ and append a shortest path from $v$ to $w$. The result is a path from $u$ to $w$ of cost $\dist(u,v)+\dist(v,w)$, and the shortest path from $u$ to $w$ cannot cost more. \qedhere
\end{proof}

This small fact is why consistent heuristics work (\cref{ch:ch04}) and why the backward search of \cref{ch:ch03} yields an exact heuristic.

\subsection{Grid graphs}
\label{sec:ch02-grids}

Most maps in this book are \textbf{occupancy grids}\index{grid!occupancy}: the workspace is cut into unit squares, and each square is either free or blocked. A grid is a graph in disguise, as \cref{def:ch02-grid} makes precise. Our coordinate convention is the mathematical one: $x$ grows to the right, $y$ grows upward, the cell $(0,0)$ sits in the lower-left corner, and we identify a cell with its center when we measure distances. The ASCII maps used by the code (\code{.} for free, \code{\#} for blocked) list the \emph{top} row first, so the code flips the row index when it reads them.

\begin{definition}[Grid graph, 4- and 8-connectivity, corner-cutting rule]
\label{def:ch02-grid}
A \textbf{grid}\index{grid} of width $W$ and height $H$ is the set of \textbf{cells} $C=\set{(x,y) : x\in\set{0,\dots,W-1},\,y\in\set{0,\dots,H-1}}$, together with a set $O\subseteq C$ of \textbf{obstacle cells}\index{grid!obstacle cell}; the cells in $C\setminus O$ are \emph{free}. The \textbf{4-connected grid graph}\index{grid!4-connected} $G_4$ has the free cells as vertices and an undirected edge of cost~$1$ between two free cells that share a side (a \emph{straight move}). The \textbf{8-connected grid graph}\index{grid!8-connected} $G_8$ has, in addition, an edge of cost~$\sqrt{2}$ between two free cells that share only a corner (a \emph{diagonal move}), subject to the \textbf{corner-cutting rule}\index{corner cutting}: the diagonal move from $(x,y)$ to $(x+d_x,\,y+d_y)$ is allowed only if the two cells $(x+d_x,\,y)$ and $(x,\,y+d_y)$ it passes between are both free.
\end{definition}

\begin{figure}[tb]
  \centering
  \inputfigure{ch02/grid-connectivity}
  \caption{Grid connectivity. (a)~The four straight moves of $G_4$ cost $1$ each. (b)~In $G_8$ the diagonal moves cost $\sqrt{2}$; the move into the obstacle is forbidden because the cell is blocked, and the two diagonals past its corners are forbidden by the corner-cutting rule. (c)~Why: a disk that slides diagonally past an obstacle corner clips it, even though the start and end cells are free.}
  \label{fig:ch02-grid-connectivity}
\end{figure}

Three choices in this definition deserve a comment. First, the diagonal cost is $\sqrt2$ and not $1$ because $\sqrt2$ is the distance between the centers of two diagonal neighbors. With costs $1$ and $\sqrt2$ the cost of a grid path is exactly the Euclidean length of the polyline through the cell centers, and the straight-line and octile heuristics of \cref{ch:ch04} never overestimate it. Second, the corner-cutting rule exists because a drone is not a point: \cref{fig:ch02-grid-connectivity}(c) shows a disk sliding diagonally past an obstacle corner and clipping it. Some code bases forbid the diagonal only when \emph{both} side cells are blocked; that rule is unsafe for a robot of positive size and this book never uses it. Third, \cref{def:ch02-grid} treats the grid as undirected: every move can be reversed. Later chapters add directed edges only when time enters (\cref{sec:ch02-time}).

Which connectivity should you use? \Cref{tab:ch02-connectivity} lists the \textbf{stretch}\index{grid!stretch factor} of each lattice, the worst-case ratio between the cost of the cheapest lattice path and the straight-line distance in an obstacle-free space. On $G_4$ the worst direction is the diagonal, where the grid path is $\sqrt2\approx1.41$ times longer than the straight line; $G_8$ brings the worst case down to $1.08$, at $22.5^\circ$ from an axis, which is why single-agent planners prefer it. The multi-agent benchmarks of \cref{ch:ch07}, following \textcite{stern2019mapf}, use $G_4$, because a unit cost per move makes ``one move per time step'' and ``one unit of cost per time step'' the same thing, which keeps conflicts simple.

\begin{table}[tb]
\centering\small
\caption{Neighborhoods on grids and lattices. The stretch is the worst-case ratio between the cost of the cheapest lattice path and the straight-line distance in an empty space, attained in the direction that the available moves imitate worst. The four values are printed by \code{worked\_example} in \code{ch02\_toolbox.py} as \code{stretch2}, \code{stretch3}, \code{stretch2\_straight} and \code{stretch3\_straight}.}
\label{tab:ch02-connectivity}
\begin{tabular}{@{}lclll@{}}
\toprule
Lattice & Moves & Move costs & Corner-cutting check & Stretch\\
\midrule
2D, 4-connected & 4 & $1$ & none & $\sqrt{2}\approx1.414$\\
2D, 8-connected & 8 & $1,\ \sqrt{2}$ & 2 side cells free & $1.082$\\
3D, 6-connected & 6 & $1$ & none & $\sqrt{3}\approx1.732$\\
3D, 26-connected & 26 & $1,\ \sqrt{2},\ \sqrt{3}$ & 2 or 6 box cells free & $1.128$\\
\bottomrule
\end{tabular}
\end{table}

\subsection{Lattices in three dimensions}
\label{sec:ch02-lattices}

A drone flies in three dimensions, and the same construction works there.

\begin{definition}[Lattice in three dimensions, 6- and 26-connectivity]
\label{def:ch02-lattice}
A \textbf{lattice}\index{lattice!3D} of size $W\times H\times D$ has the cells $(x,y,z)$ with $0\le x<W$, $0\le y<H$, $0\le z<D$, some of which are obstacle cells. A \emph{move} changes each coordinate by $-1$, $0$ or $+1$ and is not the zero move. The \textbf{6-connected lattice}\index{lattice!6-connected} allows the six moves with one non-zero component, at cost~$1$. The \textbf{26-connected lattice}\index{lattice!26-connected} allows all $26$ moves; a move with $k$ non-zero components costs $\sqrt{k}$, and the corner-cutting rule requires every cell of the axis-aligned box spanned by the two endpoints to be free. The intermediate \emph{18-connected} lattice allows the moves with at most two non-zero components.
\end{definition}

The box rule generalizes the two side cells of the plane: a move with two non-zero components spans a $2\times2$ box and checks the two other cells in it, one with three spans a $2\times2\times2$ box and checks six. In an empty 26-connected lattice the cheapest path to a cell at displacement $(a,b,c)$ with $a\ge b\ge c\ge0$ uses $c$ triple-diagonal, $b-c$ double-diagonal and $a-b$ straight moves, so it costs $(a-b)+\sqrt2\,(b-c)+\sqrt3\,c$; maximizing the ratio of this cost to $\sqrt{a^2+b^2+c^2}$ over all directions gives the stretch $1.128$ of \cref{tab:ch02-connectivity}. The price of 26-connectivity is a branching factor of $26$ instead of $6$, which multiplies the work of every search by more than four. \Cref{ch:ch16,ch:ch17} show the alternative for 3D problems: do not discretize at all and sample the free space instead.

% ---------------------------------------------------------------------
\section{Workspace and configuration space}
\label{sec:ch02-cspace}

The map lives in the \textbf{workspace}\index{workspace} $\mathcal{W}\subseteq\R^2$ or $\R^3$, the physical space in which obstacles occupy a region $\mathcal{O}\subset\mathcal{W}$. A drone is not a point of this space but a body that occupies a region, so planning a route for its center and then flying it would scrape every wall. The classical remedy, due to \textcite{lozanoperez1983spatial} and presented in every planning textbook \cite{lozanoperez1983spatial,lavalle2006planning,choset2005principles}, is to move the problem into the configuration space.

\begin{definition}[Configuration space, free space, obstacle region]
\label{def:ch02-configuration-space}
A \textbf{configuration}\index{configuration} $q$ of a robot is a tuple of parameters that fixes the position of every point of the robot; the set of all configurations is the \textbf{configuration space}\index{configuration space} $\Cspace$. Write $\mathcal{A}(q)\subset\mathcal{W}$ for the region of the workspace occupied by the robot in configuration $q$. The \textbf{obstacle region}\index{configuration space!obstacle region} and the \textbf{free space}\index{configuration space!free space} are
\begin{equation}
\Cobs=\set{q\in\Cspace : \mathcal{A}(q)\cap\mathcal{O}\neq\emptyset},\qquad
\Cfree=\Cspace\setminus\Cobs .
\label{eq:ch02-cfree}
\end{equation}
A path of configurations is \emph{collision-free} if it lies entirely in $\Cfree$.
\end{definition}

For a disk-shaped robot of radius $r$ whose orientation does not matter, the configuration is the center $\pos\in\R^2$, $\Cspace=\R^2$ and $\mathcal{A}(\pos)=\disc{\pos}{r}$, the closed disk of radius $r$ around $\pos$. The obstacle region then has a simple description, which needs one more definition.

\begin{definition}[Minkowski sum, inflation]
\label{def:ch02-minkowski-sum}
The \textbf{Minkowski sum}\index{Minkowski sum} of two sets $A,B\subseteq\R^n$ is $A\oplus B=\set{\vect{a}+\vect{b} : \vect{a}\in A,\ \vect{b}\in B}$. \textbf{Inflating}\index{obstacle inflation} a closed set $A$ by the radius $r$ means forming $A\oplus\disc{\vect{0}}{r}=\set{\vect{x} : \dist(\vect{x},A)\le r}$, the set of points within distance $r$ of $A$. Here $\dist(\vect{x},A)=\inf_{\vect{a}\in A}\norm{\vect{x}-\vect{a}}$ is the Euclidean distance from a point to a set; the same symbol $\dist$ denotes graph distance in \cref{def:ch02-shortest-path} and point-to-segment distance in \cref{def:ch02-point-segment-distance}, and which one is meant is always clear from its arguments.
\end{definition}

\begin{proposition}[Point-robot planning on the inflated map is exact]
\label[proposition]{thm:ch02-inflation}
For a disk (in 3D: a ball) robot of radius $r$ whose configuration is its center $\pos$, and a closed obstacle region $\mathcal{O}$,
\begin{equation}
\Cobs=\mathcal{O}\oplus\disc{\vect{0}}{r}.
\label{eq:ch02-inflation}
\end{equation}
Consequently $\pos\in\Cfree$ if and only if $\dist(\pos,\mathcal{O})>r$, and a path of the center is collision-free for the disk robot if and only if it is collision-free for a point robot on the map whose obstacles are inflated by $r$.
\end{proposition}
\begin{proof}
The disk $\disc{\pos}{r}$ meets $\mathcal{O}$ if and only if there is an obstacle point $\vect{o}\in\mathcal{O}$ with $\norm{\pos-\vect{o}}\le r$, that is, if and only if $\pos=\vect{o}+\vect{b}$ for some $\vect{o}\in\mathcal{O}$ and some $\vect{b}$ with $\norm{\vect{b}}\le r$. The last statement says precisely that $\pos\in\mathcal{O}\oplus\disc{\vect{0}}{r}$. The second claim applies the first to every point of the path. \qedhere
\end{proof}

\begin{figure}[tb]
  \centering
  \inputfigure{ch02/minkowski-inflation}
  \caption{Obstacle inflation. (a)~In the workspace the disk robot of radius $r$ can at best touch the obstacle. (b)~In the configuration space the obstacle grows by $r$ (note the rounded corners) and the robot shrinks to a point; the path that touches the inflated boundary corresponds to the touching disk in~(a). (c)~On a grid, every cell whose center lies within $r=1.5$ cells of an obstacle square is blocked: the two-cell obstacle drawn here blocks $16$ further cells (the text counts $5$, $9$ and $21$ for one cell at $r=0.6$, $1.0$, $1.6$).}
  \label{fig:ch02-minkowski}
\end{figure}

\Cref{fig:ch02-minkowski} shows the construction. Inflating a polygon by a disk rounds its corners, and inflating a union of obstacles inflates each of them. The same idea serves twice more: \cref{ch:ch12} turns an encounter between two disks into a point and a disk of radius $r_A+r_B$ (\cref{thm:ch02-disc-minkowski}), and in \cref{ch:ch24} the inflation radius grows with the uncertainty of a predicted position.

\paragraph{The drone as a bounding sphere.}
A quadrotor with its rotor guards is not a disk, and it rotates about its vertical axis while flying. Modeling it exactly would add the yaw angle to the configuration, giving $\Cspace=\R^3\times S^1$ and a heading-dependent obstacle region \cite{lavalle2006planning}. For a drone that is unnecessary: the body fits into a \textbf{bounding sphere}\index{bounding sphere} of radius $r$ that is the same in every heading, so the configuration is the position alone, $\Cspace=\R^3$ (or $\R^2$ at constant altitude), and \cref{thm:ch02-inflation} applies with the sphere radius. In this book every drone is such a sphere. Its radius includes a safety margin, as the following pitfall explains.

\paragraph{Inflation on a grid.}
On a grid the inflation is approximated cell by cell: a cell is added to $O$ if its center lies within $r$ of some obstacle square, which is what \code{Grid.inflate} does. Blocking a single cell and inflating by $r=0.6$, $1.0$ and $1.6$ cell widths blocks $5$, $9$ and $21$ cells respectively: a plus shape, a $3\times3$ block, and the block with a ring of $12$ further cells. The grid version is only a sampled Minkowski sum; a free cell may still contain points that are closer than $r$ to an obstacle. This is harmless for straight moves between cell centers, because obstacle centers and move endpoints both lie on the integer lattice, so for grid-aligned squares the distance to the obstacle along such a move is smallest at one of its endpoints, and the corner-cutting rule guards the diagonal moves.

\begin{pitfall}[Inflating by the body radius alone]
The radius you inflate by is a promise that the drone's center will never leave the inflated free space. Three things break that promise: the drone is larger than its nominal radius when propellers and guards are counted, the position estimate is off by the localization error, and the controller does not track the plan exactly. Use $r = r_{\text{body}} + \text{margin}$, where the margin covers the localization and tracking errors you expect, and on a grid round the result \emph{up} to the next cell width. A margin that is too small shows up as scraped walls; one that is too large closes narrow passages and makes the planner report that no path exists. \Cref{ch:ch25} measures the resulting minimum separation.
\end{pitfall}

% ---------------------------------------------------------------------
\section{Time: space-time states, paths, plans and trajectories}
\label{sec:ch02-time}

A path says \emph{where} to go; a multi-agent plan must also say \emph{when}. The book handles time in two ways. Search-based planners (Parts~II and~III) use discrete time: they fix a time step $\dt$ and let every action take exactly one step. Controllers and reactive methods (Parts~IV, VI and VII) use continuous time and describe motion by trajectories. This section defines both and the words that connect them.

\subsection{Discrete time and the time-expanded graph}
\label{sec:ch02-space-time}

\begin{definition}[Space-time state, move and wait actions]
\label{def:ch02-space-time-state}
Fix a graph $G=(V,E)$ and a time step $\dt>0$. A \textbf{space-time state}\index{space-time state} is a pair $(v,t)$ with $v\in V$ and $t\in\N=\set{0,1,2,\dots}$; it means ``the agent is at vertex $v$ at time $t\,\dt$''. Two kinds of \textbf{actions} lead from $(v,t)$ to time $t+1$: a \textbf{move}\index{move action} along an edge $(v,v')\in E$ leads to $(v',t+1)$, and the \textbf{wait action}\index{wait action} leads to $(v,t+1)$. The move costs $c(v,v')$ and the wait costs $c_{\text{wait}}>0$ for every wait made before the agent's final arrival at its goal; by the stay-at-goal convention of \cref{def:ch02-time-indexed-path} the waits at the goal after arrival are free. In the unit-cost problems of \cref{ch:ch07}, every action costs $1$. The wait cost must be positive: a free wait lets a search postpone progress forever on an unbounded horizon.
\end{definition}

Space-time states are the reason why a planner can let two agents share a corridor: the states $(v,3)$ and $(v,7)$ are different vertices of the graph below, so occupying one does not block the other.

\begin{definition}[Time-expanded graph]
\label{def:ch02-time-expanded-graph}
The \textbf{time-expanded graph}\index{time-expanded graph} of $G$ with horizon $T_{\max}$ is the directed graph with vertex set $V\times\set{0,\dots,T_{\max}}$ and, for every $t<T_{\max}$, the \textbf{wait edges} $\bigl((v,t),(v,t+1)\bigr)$ for all $v\in V$ and the \textbf{move edges} $\bigl((u,t),(v,t+1)\bigr)$ for all $(u,v)\in E$, weighted as in \cref{def:ch02-space-time-state}.
\end{definition}

\begin{figure}[tb]
  \centering
  \inputfigure{ch02/time-expanded}
  \caption{A three-vertex graph $G$ (left) and its time-expanded graph with three time layers (right). Gray arrows are wait edges, black arrows are move edges; every edge goes from layer $t$ to layer $t+1$, so the graph is acyclic. The highlighted time-indexed path waits at $a$ for one step and then moves to $b$, ending in the state $(b,2)$.}
  \label{fig:ch02-time-expanded}
\end{figure}

\Cref{fig:ch02-time-expanded} shows a base graph with three vertices and its time-expanded graph with three layers. Note what the construction costs: $\abs{V}\,(T_{\max}+1)$ vertices and $T_{\max}\,(\abs{V}+2\abs{E})$ edges for an undirected $G$ (\cref{exr:ch02-time-expanded}). Nobody stores this graph. The planner keeps $G$ and generates the successors of $(v,t)$ on demand with a function like \code{space\_time\_successors} in the toolbox, which returns the wait successor and one move successor per free neighbor, all at time $t+1$. Since all edges point forward in time, the time-expanded graph is acyclic; a search in it terminates even without a closed set, though a closed set still saves work.

\begin{definition}[Time-indexed path, arrival time, stay-at-goal convention]
\label{def:ch02-time-indexed-path}
A \textbf{time-indexed path}\index{time-indexed path} for an agent with start $s$ and goal $\gamma$ is a sequence $\pi=(v_0,v_1,\dots,v_{T_{\max}})$ with $v_0=s$, $v_{T_{\max}}=\gamma$ and, for every $t<T_{\max}$, either $v_{t+1}=v_t$ (a wait) or $(v_t,v_{t+1})\in E$ (a move); the last index $T_{\max}$ is the \textbf{horizon}\index{horizon} of the stored list. The agent occupies $v_t$ at time $t$, and by the \textbf{stay-at-goal convention}\index{stay-at-goal convention} it occupies $\gamma$ at every time $t>T_{\max}$. The \textbf{arrival time}\index{arrival time} $\cost(\pi)$ is the smallest $t$ such that $v_{t'}=\gamma$ for all $t'\ge t$; an agent that passes through $\gamma$ and leaves again has not arrived. It carries the cost operator because with unit action costs the two agree (\cref{def:ch02-travel-time}), and because the letter $T$ is reserved for horizons throughout the book. The two indices differ: $\cost(\pi)\le T_{\max}$, with equality exactly when the path ends without a trailing wait at the goal. \Cref{ch:ch07} writes the horizon of agent $i$'s stored path as $T_i$ and the plan horizon as $T=\max_i T_i$.
\end{definition}

The stay-at-goal convention matters for multi-agent planning: an agent that has landed still occupies its cell, and a later agent must route around it (\cref{ch:ch07,ch:ch08}). The arrival time is the individual cost that the multi-agent objectives of \cref{sec:ch02-costs} add up. It is computed by \code{path\_time} in the toolbox, which walks backward from the end of the sequence while the vertex equals the goal.

\subsection{Plans and trajectories}
\label{sec:ch02-plans}

\begin{definition}[Plan]
\label{def:ch02-plan}
A \textbf{plan}\index{plan} for $k$ agents is a tuple $\Pi=(\pi_1,\dots,\pi_k)$ of time-indexed paths, one per agent, all with the same time step. The position of agent $i$ at time $t$ is $v^i_t$ (\cref{ch:ch07} writes the same position as $\pi_i[t]$), extended by the stay-at-goal convention when $t$ exceeds the horizon of $\pi_i$. Which pairs of positions count as a \emph{conflict} is a modeling choice fixed in \cref{ch:ch07}; this chapter only measures how close the agents come.
\end{definition}

\begin{definition}[Trajectory]
\label{def:ch02-trajectory}
A \textbf{trajectory}\index{trajectory} is a function $\pos\colon[0,t_f]\to\Cfree$ that gives the position at every continuous time. In this book a trajectory is represented by a sequence of \textbf{waypoints}\index{waypoint} $(t_j,\pos_j,\vel_j)$ with strictly increasing times and a rule for the motion between them: constant velocity, or the double integrator of \cref{sec:ch02-motion}. A time-indexed path becomes a trajectory by straight-line motion at constant speed between consecutive vertices,
\begin{equation}
\pos\bigl((t+s)\,\dt\bigr)=(1-s)\,\pos(v_t)+s\,\pos(v_{t+1}),\qquad s\in[0,1],
\label{eq:ch02-path-to-trajectory}
\end{equation}
where $\pos(v)$ is the position of vertex $v$ in the workspace.
\end{definition}

\begin{table}[tb]
\centering\small
\caption{Path, time-indexed path, plan and trajectory. The columns say what each object fixes, which chapters produce it and which chapters consume it.}
\label{tab:ch02-path-plan-trajectory}
\begin{tabularx}{\textwidth}{@{}lYlY@{}}
\toprule
Object & Fixes & Produced by & Consumed by\\
\midrule
Path (\cref{def:ch02-path}) & where, not when & \cref{ch:ch03,ch:ch04,ch:ch05,ch:ch16} & timing, a tracking controller\\
Time-indexed path (\cref{def:ch02-time-indexed-path}) & where and when, in steps of $\dt$ & \cref{ch:ch04,ch:ch08,ch:ch09} & conflict checks, execution\\
Plan (\cref{def:ch02-plan}) & one time-indexed path per agent & \cref{ch:ch08,ch:ch09,ch:ch10,ch:ch11} & execution, the metrics of \cref{ch:ch25}\\
Trajectory (\cref{def:ch02-trajectory}) & position at every time & \cref{ch:ch13,ch:ch14,ch:ch21} & the flight controller\\
\bottomrule
\end{tabularx}
\end{table}

\Cref{tab:ch02-path-plan-trajectory} summarizes the vocabulary; keep the three words apart.

% ---------------------------------------------------------------------
\section{Cost measures}
\label{sec:ch02-costs}

\begin{definition}[Path length]
\label{def:ch02-path-length}
The \textbf{length}\index{path!length} of a path $\pi=(v_0,\dots,v_k)$ whose vertices have workspace positions $\pos(v_i)$ is
\begin{equation}
L(\pi)=\sum_{i=1}^{k}\norm{\pos(v_i)-\pos(v_{i-1})}.
\label{eq:ch02-path-length}
\end{equation}
Waiting adds nothing to the length. On a grid with the costs of \cref{def:ch02-grid} a path that contains no wait has $L(\pi)=\cost(\pi)$; a time-indexed path with $w$ waits before its arrival has $\cost(\pi)=L(\pi)+w\,c_{\text{wait}}$, so for unit costs the path cost is exactly the arrival time while $L(\pi)$ counts only the distance flown (in \cref{ex:ch02-costs}, $L(\pi_B)=2$ but $\cost(\pi_B)=3$).
\end{definition}

\begin{definition}[Travel time]
\label{def:ch02-travel-time}
The \textbf{travel time}\index{travel time} of a time-indexed path $\pi$ is its arrival time $\cost(\pi)$ of \cref{def:ch02-time-indexed-path}, measured in steps; in seconds it is $\cost(\pi)\,\dt$. It counts moves and waits alike.
\end{definition}

With unit move costs, and with the waits at the goal free as in \cref{def:ch02-space-time-state}, the accumulated cost of \cref{eq:ch02-path-cost} is exactly this travel time: a search that adds up action costs in $\gcost$ returns the arrival time, which is why the single symbol $\cost(\pi)$ serves for both. That is the convention of \textcite{stern2019mapf}, what \cref{ch:ch07} writes as $\cost(\pi_i)$, and what the notation table lists under $\sumcost$.

\begin{definition}[Makespan]
\label{def:ch02-makespan}
The \textbf{makespan}\index{makespan} of a plan $\Pi=(\pi_1,\dots,\pi_k)$ is the arrival time of the last agent,
\begin{equation}
\makespan(\Pi)=\max_{1\le i\le k} \cost(\pi_i).
\label{eq:ch02-makespan}
\end{equation}
\end{definition}

\begin{definition}[Sum of costs]
\label{def:ch02-sum-of-costs}
The \textbf{sum of costs}\index{sum of costs} of a plan $\Pi=(\pi_1,\dots,\pi_k)$ is the total travel time of all agents,
\begin{equation}
\sumcost(\Pi)=\sum_{i=1}^{k} \cost(\pi_i).
\label{eq:ch02-sumcost}
\end{equation}
\end{definition}

Later chapters write these two objectives with the macros $\makespan$ and $\sumcost$, exactly as above. They are the two standard objectives of multi-agent path finding \cite{stern2019mapf}; the sum of costs is the default in \cref{ch:ch09,ch:ch10}. They do not agree in general: a plan that minimizes one may not minimize the other (\cref{exr:ch02-costs}).

\begin{definition}[Separation, minimum separation, collision distance]
\label{def:ch02-min-separation}
For two agents with positions $\pos_i(t)$ and $\pos_j(t)$, the \textbf{separation}\index{separation} at time $t$ is $d_{ij}(t)=\norm{\pos_i(t)-\pos_j(t)}$. The \textbf{minimum separation}\index{separation!minimum} of a plan or of a set of trajectories is $d_{\min}=\min_{t}\min_{i<j} d_{ij}(t)$ over the whole execution. Two disk (or sphere) agents of radii $r_i$ and $r_j$ \textbf{collide} at time $t$ if $d_{ij}(t)<r_i+r_j$; the sum $R_{ij}=r_i+r_j$ is their \textbf{collision distance}\index{collision distance}. The touching case $d_{ij}(t)=R_{ij}$ counts as safe here and in \cref{ch:ch12,ch:ch13}, which test the strict inequality; \cref{thm:ch02-disc-minkowski} calls it an overlap only because two closed disks that touch share a boundary point.
\end{definition}

For a plan on a graph the separation can be evaluated at the integer times only (the toolbox function \code{min\_separation}) or along the straight-line motion of \cref{eq:ch02-path-to-trajectory} between consecutive steps (\code{min\_separation\_continuous}, which uses the closest-approach formula of \cref{sec:ch02-geometry} on every interval). The two answers differ, as the worked example shows.

\begin{example}[Two agents on a small grid]
\label[example]{ex:ch02-costs}
Take an empty 4-connected grid with $W=4$ and $H=3$. Agent $A$ goes from $(0,1)$ to $(3,1)$ along the middle row, agent $B$ from $(1,2)$ down to $(1,0)$. If both move without waiting, $\pi_A=\bigl((0,1),(1,1),(2,1),(3,1)\bigr)$ and $\pi_B'=\bigl((1,2),(1,1),(1,0)\bigr)$ put both agents in cell $(1,1)$ at $t=1$: the separation is $0$. Let $B$ wait one step first, $\pi_B=\bigl((1,2),(1,2),(1,1),(1,0)\bigr)$. \Cref{tab:ch02-costs-trace} traces the plan $\Pi=(\pi_A,\pi_B)$. The lengths are $L(\pi_A)=3$ and $L(\pi_B)=2$, the travel times are $\cost(\pi_A)=3$ and $\cost(\pi_B)=3$, hence $\makespan(\Pi)=3$ and $\sumcost(\Pi)=6$. At integer times the minimum separation is $1$. Between $t=1$ and $t=2$, however, $A$ moves from $(1,1)$ to $(2,1)$ while $B$ moves from $(1,2)$ to $(1,1)$; in the middle of that step, at $t=1.5$, the agents are at $(1.5,1)$ and $(1,1.5)$, only $\sqrt{0.5}\approx0.707$ apart. With a body radius of $0.3$ cells (collision distance $0.6$) the plan is safe; with a radius of $0.4$ (collision distance $0.8$) the two agents collide, since their separation $0.707$ drops below the collision distance, although they never share a cell.
\end{example}

\begin{table}[tb]
\centering\small
\caption{Trace of the plan of \cref{ex:ch02-costs}. The separation is evaluated at integer times and, in the last column, at the closest approach during the following step.}
\label{tab:ch02-costs-trace}
\begin{tabular}{@{}cccccl@{}}
\toprule
$t$ & $A$ & $B$ & $d_{AB}(t)$ & closest approach in $[t,t+1]$ & Comment\\
\midrule
0 & $(0,1)$ & $(1,2)$ & $\sqrt{2}\approx1.414$ & $1.000$ at $t=1$ & $B$ waits\\
1 & $(1,1)$ & $(1,2)$ & $1.000$ & $0.707$ at $t=1.5$ & perpendicular crossing\\
2 & $(2,1)$ & $(1,1)$ & $1.000$ & $1.000$ at $t=2$ & \\
3 & $(3,1)$ & $(1,0)$ & $\sqrt{5}\approx2.236$ & --- & both have arrived\\
\bottomrule
\end{tabular}
\end{table}

\begin{pitfall}[A conflict-free plan is not automatically collision-free]
The conflict definitions of \cref{ch:ch07} look at integer times and at swaps along an edge. \Cref{ex:ch02-costs} shows a plan without any such conflict in which two agents of radius $0.4$ cells still collide halfway through a step, their separation dropping below the collision distance. Grid planning is physically safe only if the agents are small compared with the cell (a collision distance below $\sqrt{0.5}\approx0.71$ cells for perpendicular crossings), or if you inflate the cells, add a check along the straight-line motion, or forbid such crossings as an extra conflict type. \Cref{ch:ch25} reports the minimum separation of an execution precisely so that this gap between the model and the physics becomes visible.
\end{pitfall}

% ---------------------------------------------------------------------
\section{Motion models for a drone}
\label{sec:ch02-motion}

Graph search ignores dynamics: a move is a move, whatever the speed. The local safety layer and the controller cannot afford that, because a drone at $5$~m/s needs meters to stop. The two motion models below are both \emph{linear}, which is what makes the Kalman filter of \cref{ch:ch18} and the model predictive controller of \cref{ch:ch21} tractable.

\begin{definition}[Single integrator]
\label{def:ch02-single-integrator}
The \textbf{single integrator}\index{single integrator} (point-mass model) has the state $\pos\in\R^n$, $n\in\set{2,3}$, and the velocity $\vel\in\R^n$ as its control input, limited by $\norm{\vel}\le v_{\max}$. In continuous time $\dot{\pos}=\vel$; in discrete time with step $\dt$,
\begin{equation}
\pos_{k+1}=\pos_k+\dt\,\vel_k .
\label{eq:ch02-single-integrator}
\end{equation}
\end{definition}

The single integrator assumes that the drone can change its velocity instantly. It is the model behind velocity obstacles and optimal reciprocal collision avoidance (\orca; \cref{ch:ch12,ch:ch13}), and implicitly the model of a grid planner, where one cell per step is a velocity of one cell per $\dt$.

\begin{definition}[Double integrator]
\label{def:ch02-double-integrator}
The \textbf{double integrator}\index{double integrator} has the state $\state=(\pos,\vel)\in\R^{2n}$ and the acceleration $\acc\in\R^n$ as its control input, limited by $\norm{\acc}\le a_{\max}$; the speed is limited by $\norm{\vel}\le v_{\max}$. In continuous time $\dot{\pos}=\vel$ and $\dot{\vel}=\acc$. If the acceleration is held constant during each step of length $\dt$ (a \emph{zero-order hold}), the discrete-time model is exactly
\begin{equation}
\state_{k+1}=\mat{A}\,\state_k+\mat{B}\,\acc_k,\qquad
\mat{A}=\begin{pmatrix}\mat{I}&\dt\,\mat{I}\\ \mat{0}&\mat{I}\end{pmatrix},\qquad
\mat{B}=\begin{pmatrix}\tfrac12\dt^2\,\mat{I}\\ \dt\,\mat{I}\end{pmatrix},
\label{eq:ch02-double-integrator}
\end{equation}
where $\mat{I}$ is the $n\times n$ identity.
\end{definition}

To see where $\mat{A}$ and $\mat{B}$ come from, integrate a constant acceleration $\acc$ over one step: the velocity becomes $\vel+\dt\,\acc$ and the position, which is the integral of the velocity, becomes $\pos+\dt\,\vel+\tfrac12\dt^2\acc$. The two rows of \cref{eq:ch02-double-integrator} say exactly this. The model is exact, not an approximation; the forward-Euler update $\pos_{k+1}=\pos_k+\dt\,\vel_k$ that many programs use drops the term $\tfrac12\dt^2\acc$ and is exact only when the acceleration is zero. The Kalman filter of \cref{ch:ch18} uses the same $\mat{A}$ with the acceleration replaced by noise (the \emph{constant-velocity model}), and the controller of \cref{ch:ch21} uses $\mat{A}$ and $\mat{B}$ as its prediction model.

\paragraph{Holonomic and non-holonomic.}
Not every robot can move in every direction; the distinction has a name.
\begin{definition}[Holonomic and non-holonomic systems]
\label{def:ch02-holonomic}
A robot is \textbf{holonomic}\index{holonomic} if every constraint on its motion restricts only the configurations it may occupy, so that from any configuration it can move in every direction of the configuration space. It is \textbf{non-holonomic}\index{non-holonomic} if some constraint restricts its velocity without restricting the reachable configurations.
\end{definition}
A car is the textbook non-holonomic system: it cannot slide sideways, yet it can reach any parking spot by maneuvering. Planning for such systems needs the tools of \textcite{lavalle2006planning}, which this book does not cover. A quadrotor, at the level of position, is holonomic: it can accelerate in any direction by tilting, and it can hover. That is why a drone can be abstracted as a point mass at all, while a fixed-wing aircraft, which must keep flying forward, cannot.

\paragraph{Why a quadrotor is a double integrator, and when it is not.}
A quadrotor produces one thrust force along its body axis and controls its direction by tilting. Its center of mass therefore obeys $m\ddot{\pos}=\vect{f}_{\text{thrust}}-m g\,\vect{e}_z$: given a desired acceleration $\acc$, the flight controller solves for the thrust magnitude and the attitude that produce it. Because the attitude loop of a small quadrotor settles in a fraction of a second, the position dynamics seen by a planner that replans every $\dt=0.1$~s or slower are very nearly the double integrator, with $a_{\max}$ set by the thrust-to-weight ratio and the maximum tilt. The theory of differential flatness makes this precise: any sufficiently smooth position trajectory can be tracked by a quadrotor \cite{mellinger2011minimum}. The caveats are that the acceleration cannot jump (the attitude must first rotate), and that the limits are not symmetric (gravity helps descending and hinders climbing). Choose $a_{\max}$ and $v_{\max}$ with a margin below the physical limits, and the abstraction holds.

\begin{proposition}[Stopping distance]
\label[proposition]{thm:ch02-stopping-distance}
A double integrator moving at speed $v$ that brakes with the maximal deceleration $a_{\max}$ comes to rest after the time $v/a_{\max}$ and the \textbf{stopping distance}\index{stopping distance}
\begin{equation}
d_{\text{stop}}(v)=\frac{v^2}{2\,a_{\max}} .
\label{eq:ch02-stopping-distance}
\end{equation}
Consequently a drone that must be able to stop before an obstacle at distance $d$ straight ahead may not move faster than $\sqrt{2\,a_{\max}\,d}$.
\end{proposition}
\begin{proof}
Along the direction of motion the speed is $v(t)=v-a_{\max}t$, which reaches zero at $t_s=v/a_{\max}$. The distance covered is $\int_0^{t_s} v(t)\,dt = v\,t_s-\tfrac12 a_{\max}t_s^2 = v^2/a_{\max}-v^2/(2a_{\max}) = v^2/(2a_{\max})$, which is \cref{eq:ch02-stopping-distance}. Solving $d_{\text{stop}}(v)\le d$ for $v$ gives the bound. \qedhere
\end{proof}

\begin{figure}[tb]
  \centering
  \inputfigure{ch02/double-integrator}
  \caption{A one-dimensional double integrator with $a_{\max}=1$~m/s$^2$, $v_{\max}=2$~m/s and $\dt=0.1$~s travels $10$~m. It accelerates until it reaches $v_{\max}$ at $t=2$~s, cruises, and starts braking at $t=5$~s when the remaining distance equals the stopping distance $v_{\max}^2/(2a_{\max})=2$~m (shaded). It stops exactly at $10$~m at $t=7$~s, because \cref{eq:ch02-double-integrator} is exact for piecewise-constant acceleration.}
  \label{fig:ch02-double-integrator}
\end{figure}

\Cref{fig:ch02-double-integrator} shows the formula at work. The stopping distance is the single most important number of the local safety layer: it sets the safety horizon after which an intruder must trigger avoidance (\cref{ch:ch24}), defines the admissible velocities of the dynamic window approach (\cref{ch:ch14}), and tells you how far ahead a prediction must look (\cref{ch:ch20}). \Cref{tab:ch02-motion-models} lists the models side by side.

\begin{table}[tb]
\centering\small
\caption{Motion models used in the book. The graph model has no dynamics at all; the two integrator models are linear; the full quadrotor model appears only inside a physics simulator.}
\label{tab:ch02-motion-models}
\begin{tabularx}{\textwidth}{@{}lllYY@{}}
\toprule
Model & State & Input & Limits & Used in\\
\midrule
Graph move & vertex $v$ & move or wait & one edge per step & Parts~II and~III\\
Single integrator & $\pos$ & $\vel$ & $\norm{\vel}\le v_{\max}$ & \cref{ch:ch12,ch:ch13,ch:ch15,ch:ch16,ch:ch17}\\
Double integrator & $(\pos,\vel)$ & $\acc$ & $\norm{\acc}\le a_{\max}$, $\norm{\vel}\le v_{\max}$ & \cref{ch:ch14,ch:ch18,ch:ch21,ch:ch22}\\
Unicycle (ground robot) & $(\pos,\theta)$ & speed, turn rate & non-holonomic & contrast only\\
Full quadrotor & $\pos,\vel$, attitude, rates & thrust, torques & motor limits & simulation only\\
\bottomrule
\end{tabularx}
\end{table}

\begin{pitfall}[Clipping breaks linearity]
The limits $\norm{\acc}\le a_{\max}$ and $\norm{\vel}\le v_{\max}$ are not part of the linear model \cref{eq:ch02-double-integrator}; they are constraints on top of it. Code that first applies $\mat{A}$ and $\mat{B}$ and then clips the velocity has silently changed the position update too, because the position moved with the unclipped velocity for half a step. The toolbox function \code{double\_integrator\_step} avoids the inconsistency by clipping the input, computing the new velocity, and then advancing the position with the average of the old and the new velocity, which coincides with $\mat{A}\state+\mat{B}\acc$ whenever no limit is active. Model predictive control (\cref{ch:ch21}) handles the limits properly, as constraints of an optimization problem.
\end{pitfall}

% ---------------------------------------------------------------------
\section{A geometry toolkit}
\label{sec:ch02-geometry}

The local safety layer of the book is built from a handful of geometric primitives. They are collected here with their formulas and proofs; the toolbox implements each of them on plain tuples of floats. The vector algebra they rest on---norms, dot and cross products, projections---is in \cref{sec:appB-linalg}.

\begin{definition}[Vectors, dot product, 2D cross product]
\label{def:ch02-dot-cross}
For $\vect{a},\vect{b}\in\R^n$ the \textbf{dot product}\index{dot product} is $\vect{a}\cdot\vect{b}=\sum_i a_ib_i=\norm{\vect{a}}\,\norm{\vect{b}}\cos\phi$, where $\phi$ is the angle between the vectors and $\norm{\vect{a}}=\sqrt{\vect{a}\cdot\vect{a}}$. The \textbf{projection} of $\vect{p}$ onto a unit vector $\vect{u}$ is $(\vect{p}\cdot\vect{u})\,\vect{u}$, and $\vect{p}-(\vect{p}\cdot\vect{u})\,\vect{u}$ is the part of $\vect{p}$ perpendicular to $\vect{u}$. In the plane, the \textbf{2D cross product}\index{cross product!2D} $\vect{a}\times\vect{b}=a_xb_y-a_yb_x=\norm{\vect{a}}\,\norm{\vect{b}}\sin\phi$ is the signed area of the parallelogram spanned by the two vectors; it is positive when $\vect{b}$ lies counterclockwise from $\vect{a}$ (a left turn). The vector $\vect{a}^{\perp}=(-a_y,a_x)$ is $\vect{a}$ rotated by $90^\circ$ counterclockwise. In three dimensions $\vect{a}\times\vect{b}$ is the vector perpendicular to both with length $\norm{\vect{a}}\,\norm{\vect{b}}\,\abs{\sin\phi}$.
\end{definition}

\begin{definition}[Distance from a point to a segment]
\label{def:ch02-point-segment-distance}
Let $\vect{a}$ and $\vect{b}$ be the endpoints of a segment and $\vect{p}$ a point. With
\begin{equation}
s=\min\Bigl(1,\ \max\Bigl(0,\ \frac{(\vect{p}-\vect{a})\cdot(\vect{b}-\vect{a})}{\norm{\vect{b}-\vect{a}}^2}\Bigr)\Bigr),
\qquad \vect{q}=\vect{a}+s\,(\vect{b}-\vect{a}),
\label{eq:ch02-point-segment}
\end{equation}
the point $\vect{q}$ is the \textbf{closest point}\index{point--segment distance} of the segment to $\vect{p}$ and $\dist(\vect{p},\overline{\vect{a}\vect{b}})=\norm{\vect{p}-\vect{q}}$. When $\vect{a}=\vect{b}$ we set $s=0$.
\end{definition}

The unclamped fraction in \cref{eq:ch02-point-segment} is the projection of $\vect{p}$ onto the line through $\vect{a}$ and $\vect{b}$, expressed as a multiple of $\vect{b}-\vect{a}$. Clamping it to $[0,1]$ handles the three cases: the foot of the perpendicular lies inside the segment ($0<s<1$), or before $\vect{a}$ ($s=0$), or beyond $\vect{b}$ ($s=1$). This one function is the collision checker of the book: a disk obstacle of radius $r$ blocks a straight motion from $\vect{a}$ to $\vect{b}$ if and only if the distance from its center to the segment is at most $r$, and the clearance of a candidate motion in \cref{ch:ch14} is the smallest such distance over all obstacles.

\begin{proposition}[Ray--circle and segment--circle intersection]
\label[proposition]{thm:ch02-ray-circle}
Let $\vect{o}+t\,\vect{d}$, $t\ge0$, be a ray and $\norm{\vect{x}-\vect{c}}=r$ a circle, and put $\vect{f}=\vect{o}-\vect{c}$. The ray meets the circle where
\begin{equation}
\begin{gathered}
(\vect{d}\cdot\vect{d})\,t^2+2\,(\vect{f}\cdot\vect{d})\,t+(\vect{f}\cdot\vect{f}-r^2)=0 ,\\
t_{1,2}=\frac{-(\vect{f}\cdot\vect{d})\mp\sqrt{\Delta}}{\vect{d}\cdot\vect{d}},
\qquad \Delta=(\vect{f}\cdot\vect{d})^2-(\vect{d}\cdot\vect{d})(\vect{f}\cdot\vect{f}-r^2).
\end{gathered}
\label{eq:ch02-ray-circle}
\end{equation}
If $\Delta<0$ the ray misses the circle; otherwise the first contact is at the smaller root $t_1$, provided $t_1\ge0$. If $\vect{f}\cdot\vect{f}\le r^2$ the ray starts inside the disk and the contact time is $0$. The segment from $\vect{a}$ to $\vect{b}$ meets the disk $\disc{\vect{c}}{r}$ if and only if $\dist(\vect{c},\overline{\vect{a}\vect{b}})\le r$; it crosses the circle at the roots of \cref{eq:ch02-ray-circle} with $\vect{o}=\vect{a}$, $\vect{d}=\vect{b}-\vect{a}$ that lie in $[0,1]$.
\end{proposition}
\begin{proof}
Substitute $\vect{x}=\vect{o}+t\vect{d}$ into $\norm{\vect{x}-\vect{c}}^2=r^2$ and expand $\norm{\vect{f}+t\vect{d}}^2$; the quadratic formula gives the roots. The segment statement restates the definition of the distance to a segment. \qedhere
\end{proof}

\begin{figure}[tb]
  \centering
  \inputfigure{ch02/segment-circle}
  \caption{(a)~The segment $\overline{\vect{a}\vect{b}}$ meets the disk around $\vect{c}$ because the closest point $\vect{q}$ is at the point-to-segment distance $d_1=\dist(\vect{c},\overline{\vect{a}\vect{b}})\le r$; the ray $\vect{a}+t(\vect{b}-\vect{a})$ crosses the circle at $t_1$ and $t_2$. The disk around $\vect{c}'$ is missed because its distance $d_2$ to the segment exceeds $r$. (b)~Here $d=\norm{\vect{p}-\vect{c}}$ is the distance to the \emph{center}, the quantity of \cref{thm:ch02-tangents}. The two tangents from $\vect{p}$ to a circle touch it at $\vect{t}_{\pm}$; the radius is perpendicular to the tangent, so $\cos\alpha=r/d$ and $\sin\theta=r/d$. Every ray from $\vect{p}$ inside the shaded cone hits the disk; this cone is the velocity obstacle of \cref{ch:ch12}.}
  \label{fig:ch02-segment-circle}
\end{figure}

\begin{proposition}[Minkowski sum of two disks]
\label[proposition]{thm:ch02-disc-minkowski}
$\disc{\vect{c}_1}{r_1}\oplus\disc{\vect{c}_2}{r_2}=\disc{\vect{c}_1+\vect{c}_2}{r_1+r_2}$. In particular, two disks $\disc{\pos_A}{r_A}$ and $\disc{\pos_B}{r_B}$ overlap if and only if $\norm{\pos_B-\pos_A}\le r_A+r_B$, i.e., if and only if the point $\pos_B-\pos_A$ lies in $\disc{\vect{0}}{r_A+r_B}$.
\end{proposition}
\begin{proof}
If $\vect{x}=\vect{a}+\vect{b}$ with $\norm{\vect{a}-\vect{c}_1}\le r_1$ and $\norm{\vect{b}-\vect{c}_2}\le r_2$, the triangle inequality gives $\norm{\vect{x}-\vect{c}_1-\vect{c}_2}\le r_1+r_2$. Conversely, assume $r_1+r_2>0$ (if both radii vanish the two sides are the single point $\vect{c}_1+\vect{c}_2$ and the claim is trivial); if $\vect{w}=\vect{x}-\vect{c}_1-\vect{c}_2$ has $\norm{\vect{w}}\le r_1+r_2$, take $\vect{a}=\vect{c}_1+\frac{r_1}{r_1+r_2}\vect{w}$ and $\vect{b}=\vect{c}_2+\frac{r_2}{r_1+r_2}\vect{w}$; then $\vect{a}+\vect{b}=\vect{x}$ and both points lie in their disks. The overlap statement is \cref{thm:ch02-inflation} applied with $\mathcal{O}=\disc{\pos_B}{r_B}$ and the robot $\disc{\pos_A}{r_A}$. \qedhere
\end{proof}

This is the \emph{relative frame} that \cref{ch:ch12,ch:ch13} use throughout: put drone $A$ at the origin as a point, and replace drone $B$ by a disk of radius $r_A+r_B$ at the relative position $\pos_B-\pos_A$. The next primitive describes the directions in which that disk is visible from the origin.

\begin{proposition}[Tangent lines from a point to a circle]
\label[proposition]{thm:ch02-tangents}
Let $d=\norm{\vect{p}-\vect{c}}>r$ and $\vect{u}=(\vect{p}-\vect{c})/d$. The two tangents from $\vect{p}$ to the circle $\norm{\vect{x}-\vect{c}}=r$ touch it at
\begin{equation}
\vect{t}_{\pm}=\vect{c}+r\,\bigl(\cos\alpha\,\vect{u}\pm\sin\alpha\,\vect{u}^{\perp}\bigr),
\qquad \cos\alpha=\frac{r}{d},\quad \sin\alpha=\sqrt{1-\frac{r^2}{d^2}} .
\label{eq:ch02-tangent-points}
\end{equation}
Both tangents have length $\ell=\sqrt{d^2-r^2}$ and make the half-angle $\theta=\arcsin(r/d)$ with the line from $\vect{p}$ to $\vect{c}$; a ray from $\vect{p}$ meets the disk if and only if its direction is within $\theta$ of $\vect{c}-\vect{p}$.
\end{proposition}
\begin{proof}
A radius is perpendicular to the tangent at its endpoint, so $\vect{p}$, $\vect{c}$ and $\vect{t}_{\pm}$ form a right triangle with the right angle at $\vect{t}_{\pm}$, hypotenuse $d$ and legs $r$ and $\ell$. Hence $\cos\alpha=r/d$ for the angle $\alpha$ at $\vect{c}$, $\ell=\sqrt{d^2-r^2}$, and $\sin\theta=r/d$ for the angle $\theta$ at $\vect{p}$. Rotating $\vect{u}$ by $\pm\alpha$ and scaling by $r$ gives the two tangent points relative to $\vect{c}$. The rays that hit the disk are exactly those between the two tangents. \qedhere
\end{proof}

\Cref{fig:ch02-segment-circle}(b) shows the geometry for $\vect{p}=(0,0)$, $\vect{c}=(4,0)$ and $r=2$: the tangent points are $(3,\pm\sqrt3)$, the tangent length is $\sqrt{12}$ and the half-angle is $30^\circ$ (\cref{exr:ch02-tangents}). When $\vect{p}$ approaches the circle, $\theta$ grows toward $90^\circ$ and the cone opens into a half-plane; when $\vect{p}$ is inside, there are no tangents, and the toolbox returns \code{None}.

\begin{definition}[Closest approach and time to collision]
\label{def:ch02-closest-approach}
Let two objects move with constant velocities, $\pos_A(t)=\pos_A+t\,\vel_A$ and $\pos_B(t)=\pos_B+t\,\vel_B$, and write
\begin{equation}
\pos_{\mathrm{rel}}=\pos_B-\pos_A,\qquad \vel_{\mathrm{rel}}=\vel_A-\vel_B
\label{eq:ch02-relative}
\end{equation}
for the \textbf{relative position}\index{relative position} of $B$ with respect to $A$ and the \textbf{relative velocity}\index{relative velocity} of $A$ with respect to $B$. These are the signs that \cref{ch:ch12,ch:ch13} use to build the velocity obstacle, and this book uses no others. Their distance is $D(t)=\norm{\pos_{\mathrm{rel}}-t\,\vel_{\mathrm{rel}}}$. The \textbf{time of closest approach}\index{closest approach!time of} is $t^*=\argmin_{t\ge0}D(t)$, the \textbf{minimum distance} is $d_{\min}=D(t^*)$, and, for a collision distance $R$, the \textbf{time to collision}\index{time to collision} is $t_c=\min\set{t\ge0 : D(t)\le R}$, or $\infty$ if no such time exists (the toolbox function \code{time\_to\_collision} returns \code{None} in that case). With a horizon $\ttc$ the minimization is restricted to $t\in[0,\ttc]$.
\end{definition}

\begin{proposition}[Closest approach of two constant-velocity objects]
\label[proposition]{thm:ch02-closest-approach}
If $\vel_{\mathrm{rel}}\neq\vect{0}$ then
\begin{equation}
t^*=\max\Bigl(0,\ \frac{\pos_{\mathrm{rel}}\cdot\vel_{\mathrm{rel}}}{\vel_{\mathrm{rel}}\cdot\vel_{\mathrm{rel}}}\Bigr),
\qquad
d_{\min}=\norm{\pos_{\mathrm{rel}}-t^*\vel_{\mathrm{rel}}},
\label{eq:ch02-tca}
\end{equation}
and when the maximum is not clamped, $d_{\min}=\abs{\pos_{\mathrm{rel}}\times\vel_{\mathrm{rel}}}/\norm{\vel_{\mathrm{rel}}}$ in the plane (in 3D, $\norm{\pos_{\mathrm{rel}}\times\vel_{\mathrm{rel}}}/\norm{\vel_{\mathrm{rel}}}$). The objects come within $R$ of each other if and only if $d_{\min}\le R$, and then
\begin{equation}
t_c=\max\Bigl(0,\ t^*-\frac{\sqrt{R^2-d_{\min}^2}}{\norm{\vel_{\mathrm{rel}}}}\Bigr).
\label{eq:ch02-ttc}
\end{equation}
If $\vel_{\mathrm{rel}}=\vect{0}$ the distance is constant, $t^*=0$ and $t_c$ is $0$ if $\norm{\pos_{\mathrm{rel}}}\le R$ and $\infty$ otherwise.
\end{proposition}
\begin{proof}
$D(t)^2=\norm{\pos_{\mathrm{rel}}}^2-2t\,(\pos_{\mathrm{rel}}\cdot\vel_{\mathrm{rel}})+t^2\norm{\vel_{\mathrm{rel}}}^2$ is a parabola in $t$ with positive leading coefficient; its derivative $-2(\pos_{\mathrm{rel}}\cdot\vel_{\mathrm{rel}})+2t\norm{\vel_{\mathrm{rel}}}^2$ vanishes at $t=(\pos_{\mathrm{rel}}\cdot\vel_{\mathrm{rel}})/\norm{\vel_{\mathrm{rel}}}^2$. If this value is negative, the objects are already separating and the closest approach on $t\ge0$ is at $t=0$, which the maximum expresses. At the unclamped minimum, $D^2=\norm{\pos_{\mathrm{rel}}}^2-(\pos_{\mathrm{rel}}\cdot\vel_{\mathrm{rel}})^2/\norm{\vel_{\mathrm{rel}}}^2=\bigl(\norm{\pos_{\mathrm{rel}}}^2\norm{\vel_{\mathrm{rel}}}^2-(\pos_{\mathrm{rel}}\cdot\vel_{\mathrm{rel}})^2\bigr)/\norm{\vel_{\mathrm{rel}}}^2=(\pos_{\mathrm{rel}}\times\vel_{\mathrm{rel}})^2/\norm{\vel_{\mathrm{rel}}}^2$ by Lagrange's identity. For the time to collision, $D(t)^2=R^2$ is the quadratic of \cref{thm:ch02-ray-circle} with $\vect{o}=\pos_{\mathrm{rel}}$, $\vect{d}=-\vel_{\mathrm{rel}}$, $\vect{c}=\vect{0}$; it has real roots exactly when $d_{\min}\le R$, symmetric about $t^*$ at distance $\sqrt{R^2-d_{\min}^2}/\norm{\vel_{\mathrm{rel}}}$, and the smaller root is the first contact, clamped at $0$ when the disks already overlap. \qedhere
\end{proof}

\begin{figure}[tb]
  \centering
  \inputfigure{ch02/closest-approach}
  \caption{Closest approach in the relative frame of \cref{ex:ch02-closest-approach}. Drone $A$ sits at the origin as a point; drone $B$ starts at $\pos_{\mathrm{rel}}=(4,2)$ and closes on $A$ along $-\vel_{\mathrm{rel}}=(-1,-1)$, the relative velocity being $\vel_{\mathrm{rel}}=\vel_A-\vel_B=(1,1)$. The perpendicular from the origin to the line of motion gives $d_{\min}=\sqrt2$ at $t^*=3$. With collision distance $R=1.5$ the relative motion enters the disk at $t_c\approx2.65$.}
  \label{fig:ch02-closest-approach}
\end{figure}

\begin{example}[Two drones on crossing courses]
\label[example]{ex:ch02-closest-approach}
Drone $A$ is at the origin and flies along the $x$-axis at $1$~m/s; drone $B$ is at $(4,2)$ and flies straight down at $1$~m/s. Then $\pos_{\mathrm{rel}}=(4,2)$ and $\vel_{\mathrm{rel}}=\vel_A-\vel_B=(1,0)-(0,-1)=(1,1)$. From \cref{eq:ch02-tca}, $t^*=(\pos_{\mathrm{rel}}\cdot\vel_{\mathrm{rel}})/(\vel_{\mathrm{rel}}\cdot\vel_{\mathrm{rel}})=6/2=3$~s and $d_{\min}=\abs{\pos_{\mathrm{rel}}\times\vel_{\mathrm{rel}}}/\norm{\vel_{\mathrm{rel}}}=\abs{4\cdot1-2\cdot1}/\sqrt2=2/\sqrt2=\sqrt2\approx1.414$~m. If both drones have radius $0.75$~m, the collision distance is $R=1.5>\sqrt2$, and \cref{eq:ch02-ttc} gives the first contact at $t_c=3-\sqrt{2.25-2}/\sqrt2\approx2.646$~s. If both have radius $0.5$~m then $R=1<\sqrt2$ and they pass each other safely. The toolbox functions \code{time\_of\_closest\_approach} and \code{time\_to\_collision} return exactly these numbers; \cref{fig:ch02-closest-approach} draws the situation.
\end{example}

Two remarks close the toolkit. The horizon $\ttc$ of \cref{def:ch02-closest-approach} is the form in which the local safety layer of \cref{ch:ch24} asks the question. And everything here carries over to three dimensions, except that the cross product becomes a vector (use its norm in \cref{eq:ch02-tca}) and the two tangent lines become a cone with the same half-angle $\theta=\arcsin(r/d)$.

% ---------------------------------------------------------------------
\section{Uncertainty: random vectors, Gaussians and Bayes' rule}
\label{sec:ch02-uncertainty}

The tracker of \cref{ch:ch18} does not hand the planner a position; it hands over an estimate \emph{and} a covariance that says how much to trust it. \Cref{ch:appB} owns the mathematics of that object for the whole book: \cref{sec:appB-probability} defines the mean and covariance of a random vector (\cref{def:appB-covariance}), proves the affine rule $\Cov(\mat{A}\state+\vect{b})=\mat{A}\mat{\Sigma}\mat{A}\T$, which needs no Gaussian assumption (\cref{thm:appB-cov-rules}), gives the multivariate Gaussian density and its Mahalanobis exponent (\cref{def:appB-gaussian}), and shows that Gaussians survive affine maps and the addition of independent Gaussian noise (\cref{thm:appB-affine}). Those last two closure properties are the whole reason the Kalman filter has a closed form: pushing a Gaussian through the linear motion model of \cref{def:ch02-double-integrator} and adding noise gives another Gaussian, so the filter has only a mean and a covariance to update (\cref{ch:ch18}). This section assumes all of that and spends its space on the use \cref{ch:appB} does not make of it: reading a covariance as a picture of an intruder's uncertainty, the form in which \cref{ch:ch20} hands a prediction to \cref{ch:ch24}. A covariance matrix is hard to read as numbers; it is easy to read as an ellipse.

\begin{definition}[Covariance ellipse]
\label{def:ch02-covariance-ellipse}
The \textbf{$k\sigma$ covariance ellipse}\index{covariance ellipse} (in 3D: ellipsoid) of $\Normal(\vect{\mu},\mat{\Sigma})$ is the set of points at Mahalanobis distance $k$ from the mean,
\begin{equation}
\mathcal{E}_k=\set{\state : (\state-\vect{\mu})\T\mat{\Sigma}\inv(\state-\vect{\mu})=k^2}.
\label{eq:ch02-ellipse}
\end{equation}
\end{definition}

\begin{proposition}[Axes of the covariance ellipse]
\label[proposition]{thm:ch02-ellipse-axes}
Let $\mat{\Sigma}=\mat{E}\mat{\Lambda}\mat{E}\T$ be the eigen-decomposition of a $2\times2$ covariance, with orthonormal eigenvectors $\vect{e}_1,\vect{e}_2$ as the columns of $\mat{E}$ and eigenvalues $\lambda_1\ge\lambda_2>0$. Then $\mathcal{E}_k$ is the ellipse centerd at $\vect{\mu}$ with semi-axis $k\sqrt{\lambda_1}$ along $\vect{e}_1$ and $k\sqrt{\lambda_2}$ along $\vect{e}_2$,
\begin{equation}
\state(\varphi)=\vect{\mu}+k\bigl(\sqrt{\lambda_1}\cos\varphi\;\vect{e}_1+\sqrt{\lambda_2}\sin\varphi\;\vect{e}_2\bigr),\qquad \varphi\in[0,2\pi),
\label{eq:ch02-ellipse-param}
\end{equation}
and the major axis makes the angle $\operatorname{atan2}(e_{1,y},e_{1,x})$ with the $x$-axis, with $\vect{e}_1$ normalized so that $e_{1,x}\ge0$ (an eigenvector has no canonical sign, and this is the choice \code{covariance\_ellipse} makes). In two dimensions the probability mass inside $\mathcal{E}_k$ is $1-e^{-k^2/2}$: $39.3\%$ for $k=1$, $86.5\%$ for $k=2$ and $98.9\%$ for $k=3$.
\end{proposition}
\begin{proof}
Substitute $\vect{y}=\mat{E}\T(\state-\vect{\mu})$. Since $\mat{E}$ is orthogonal, $\mat{\Sigma}\inv=\mat{E}\mat{\Lambda}\inv\mat{E}\T$ and the quadratic form becomes $y_1^2/\lambda_1+y_2^2/\lambda_2=k^2$, an axis-aligned ellipse with semi-axes $k\sqrt{\lambda_i}$, which \cref{eq:ch02-ellipse-param} parametrizes; $\state=\vect{\mu}+\mat{E}\vect{y}$ rotates it back onto the eigenvectors. For the mass, note that $\vect{y}$ scaled by $\mat{\Lambda}^{-1/2}$ is a standard Gaussian in $\R^2$, whose squared norm is $\chi^2$-distributed with two degrees of freedom (\cref{thm:appB-chi2}), i.e., exponential with mean $2$; hence $\Prob(\text{inside }\mathcal{E}_k)=1-e^{-k^2/2}$. \qedhere
\end{proof}

To draw the ellipse, call \code{numpy.linalg.eigh} on $\mat{\Sigma}$ (its eigenvalues come in ascending order, so reverse them), take the square roots and evaluate \cref{eq:ch02-ellipse-param} at, say, $72$ angles; the toolbox does this in \code{covariance\_ellipse} and \code{ellipse\_points}.

\begin{example}[Reading a covariance matrix]
\label[example]{ex:ch02-ellipse}
Let $\vect{\mu}=(2,1)$ and $\mat{\Sigma}=\begin{psmallmatrix}2&1.2\\1.2&1\end{psmallmatrix}$, so $\sigma_x=\sqrt2$, $\sigma_y=1$ and $\rho=1.2/\sqrt2\approx0.85$. The eigenvalues solve $\lambda^2-3\lambda+0.56=0$, giving $\lambda_1=2.8$ and $\lambda_2=0.2$, with $\vect{e}_1\propto(3,2)$. The $1\sigma$ ellipse therefore has semi-axes $\sqrt{2.8}\approx1.673$ and $\sqrt{0.2}\approx0.447$, and its major axis is tilted by $\arctan(2/3)\approx33.7^\circ$. \Cref{fig:ch02-gaussian} shows $400$ samples drawn with a fixed seed: $39.8\%$ of them fall inside the $1\sigma$ ellipse and $85.8\%$ inside the $2\sigma$ ellipse, close to the theoretical $39.3\%$ and $86.5\%$.
\end{example}

\begin{figure}[tb]
  \centering
  \inputfigure{ch02/gaussian-ellipse}
  \caption{Samples from the Gaussian of \cref{ex:ch02-ellipse} with its $1\sigma$ and $2\sigma$ covariance ellipses. The arrows are the semi-axes $\sqrt{\lambda_i}\,\vect{e}_i$ of the $1\sigma$ ellipse; the strong positive correlation shows as the tilt. Only about $39\%$ of the samples lie inside the $1\sigma$ ellipse.}
  \label{fig:ch02-gaussian}
\end{figure}

\begin{pitfall}[The 68--95 rule is one-dimensional]
In one dimension $68.3\%$ of the mass of a Gaussian lies within $1\sigma$ and $95.4\%$ within $2\sigma$. In two dimensions the $1\sigma$ ellipse holds only $39.3\%$ and the $2\sigma$ ellipse $86.5\%$; in three dimensions the figures drop to $19.9\%$ and $73.9\%$. If you inflate a predicted obstacle by ``one sigma'' to be safe, you are covering less than half of the cases. Choose $k$ from the $\chi^2$ distribution with the right number of degrees of freedom for the confidence you need; \cref{ch:ch24} does this when it turns predictions into constraints.
\end{pitfall}

\paragraph{Bayes' rule.}
Everything in Part~VI is an application of one formula, posterior $\propto$ likelihood $\times$ prior, stated as \cref{eq:appB-bayes} and turned into the predict--update recursion of a filter in \cref{eq:appB-predict,eq:appB-update}. Tracking applies it again and again: today's posterior, pushed through the motion model, becomes tomorrow's prior. With a Gaussian prior, a linear motion model and a linear measurement with Gaussian noise the posterior is again Gaussian and the rule reduces to a few matrix operations, which is the Kalman filter of \cref{ch:ch18}; \cref{ch:ch19} covers what to do when the models are not linear. The standard reference for all of this is \textcite{thrun2005probabilistic}.

% ---------------------------------------------------------------------
\section{Computation: complexity, priority queues and hashing}
\label{sec:ch02-computation}

Asymptotic notation, the mechanics of a binary heap and the cost model of a hash table are stated once, with their proofs, in \cref{sec:appB-complexity}: $f(n)=\bigO{g(n)}$ means $f(n)\le c\,g(n)$ beyond some $n_0$ (\cref{def:appB-bigO}), a binary heap inserts and extracts the minimum in $\bigO{\log n}$ and peeks at it in $\bigO{1}$, a hash table inserts, looks up and deletes in constant expected time, and \cref{tab:appB-heap} collects the cost of every container the book's code uses. What a \emph{search} adds is the reading. Its running time is expressed in the sizes of \emph{its} input---the numbers of vertices $\abs{V}$ and edges $\abs{E}$, the branching factor $b$, the horizon $T_{\max}$, the number of agents $k$---and which of them a bound is read in matters more than the bound itself: the exponential worst case of \cref{ch:ch09,ch:ch10} is in $k$, not in $\abs{V}$. Big-O also hides constants, and in \python they are large: a loop iteration costs about a hundred nanoseconds, while NumPy does the same work per element roughly a hundred times faster, because the loop runs in C. Two of the operations in \cref{tab:appB-heap} deserve a closer look, because getting them wrong turns an $\bigO{n\log n}$ search into an $\bigO{n^2}$ one: the missing decrease-key, and the hashing of a state.

\paragraph{The lazy-deletion idiom.}
Dijkstra's algorithm and \astar (\cref{ch:ch03,ch:ch04}) must lower the key of a vertex whenever they find a cheaper path to it, and Lifelong Planning \astar (\lpastar) and \dstarlite (\cref{ch:ch05}) must also remove vertices from the queue. \python's \code{heapq}, the binary heap of \cref{sec:appB-complexity} over a plain list, offers neither operation, and searching the list for the entry would cost $\bigO{n}$. The idiom used throughout this book, called \textbf{lazy deletion}\index{lazy deletion}\index{priority queue!lazy deletion}, avoids both: never touch an entry that is already in the heap; when a key improves, push a \emph{second} entry with the better key; when popping, skip any entry whose key is no longer the item's best known key. \Cref{alg:ch02-lazypq} states the two operations; the class \code{LazyPQ} of \code{code/ch02\_toolbox.py} implements them.

\begin{algorithm}[htb]
\caption{Priority queue with lazy deletion}
\label{alg:ch02-lazypq}
\KwData{binary heap $H$ of triples $(\key,\text{counter},\text{item})$; dictionary $\mathit{best}$ mapping items to their best pushed key; counter $c$}
\Fn{\ToolboxLazyPush{$\key$, $\text{item}$}}{
  \If{$\text{item}\in\mathit{best}$ \KwAnd $\mathit{best}[\text{item}]\le\key$}{\KwRet{}\tcp*{no improvement: ignore}}
  $\mathit{best}[\text{item}]\gets\key$; $c\gets c+1$\label{alg:ch02-lazypq:best}\;
  push $(\key,c,\text{item})$ onto $H$\tcp*{duplicates are allowed}
}
\Fn{\ToolboxLazyPop{}}{
  \While{$H$ is not empty}{
    $(\key,\cdot,\text{item})\gets$ pop the minimum of $H$\label{alg:ch02-lazypq:pop}\;
    \If{$\text{item}\in\mathit{best}$ \KwAnd $\mathit{best}[\text{item}]=\key$}{
      delete $\mathit{best}[\text{item}]$; \KwRet{$(\key,\text{item})$}\;
    }
    \tcp{otherwise the entry is stale: a better key was pushed later}
  }
  \KwRet{empty}\;
}
\end{algorithm}

Line~\ref{alg:ch02-lazypq:best} of \cref{alg:ch02-lazypq} records the best key of each item, which is what makes a stale entry recognizable at line~\ref{alg:ch02-lazypq:pop}: an entry is stale exactly when its key differs from the recorded best key. Each push costs $\bigO{\log n}$, each pop costs $\bigO{\log n}$ plus one extra pop per stale entry it skips; since every entry is popped at most once, the total cost of a search with $P$ pushes is $\bigO{P\log P}$. A search that relaxes every edge once pushes at most $\abs{E}$ entries, so the heap can hold $\abs{E}$ rather than $\abs{V}$ entries; because $\log\abs{E}\le2\log\abs{V}$, the asymptotic bound of the search is unchanged. Two variants appear in the literature: comparing the popped key with the vertex's current $\gcost$ instead of keeping a separate dictionary, which is equivalent to the version above, and skipping a popped vertex that is already in a closed set, which never re-opens a node while \cref{alg:ch02-lazypq} does. The two coincide whenever the heuristic is consistent, and \cref{ch:ch04} is where reopening is analyzed. Where \cref{alg:ch02-lazypq} returns ``empty'', \code{LazyPQ.pop} raises \code{IndexError}, so a search loop tests \code{len(pq)} before popping. The toolbox version also counts pushes and stale pops so that you can watch the overhead in the exercises.

\paragraph{Encoding a state.}
The dictionaries of a search ($\gcost$-values, parents, the closed set, the best keys of the queue) need hashable keys, and \cref{sec:appB-complexity} explains why floating-point keys are legal but dangerous. The book's convention follows from that warning: a state is always a tuple of integers, \code{(x, y)} for a grid cell, \code{((x, y), t)} for a space-time state, \code{(x, y, z)} for a lattice cell. Such tuples hash in constant time, compare quickly and cost about a hundred bytes each, and---unlike a rounded float---two physically identical states always produce the same key, so the closed set of \cref{ch:ch03,ch:ch04} cannot silently stop working. Continuous states, as in \cref{ch:ch16,ch:ch17}, are \emph{not} hashed at all; they live in arrays and are searched with a KD-tree.

\paragraph{Why dictionaries are enough.}
The largest instances in this book are grids of about $100\times100$ cells with horizons of about $100$ steps, i.e., $10^6$ space-time states, and a \python dictionary holding $10^6$ tuple keys occupies about $150$~MB and is filled in a few seconds. A search touches far fewer states than that, so no chapter uses a specialized data structure for its closed set; \cref{ch:ch25} is where faster encodings are discussed.

% ---------------------------------------------------------------------
\section{The Python stack and the code of the book}
\label{sec:ch02-python}

\begin{table}[tb]
\centering\small
\caption{The software stack. Only the first two rows are required to run the code of the book.}
\label{tab:ch02-stack}
\begin{tabularx}{\textwidth}{@{}llY@{}}
\toprule
Package & Status & Used for\\
\midrule
\python $\ge3.10$ & required & all code; tuples, dictionaries, \code{heapq}, \code{math}\\
NumPy & required & arrays, linear algebra (\code{eigh}, \code{solve}), seeded random numbers\\
SciPy (\code{scipy.optimize}) & Part~VII only & quadratic programs and \code{milp} in \cref{ch:ch21,ch:ch22}\\
Matplotlib & optional & your own plots and animations; the book's figures are TikZ\\
NetworkX \cite{hagberg2008networkx} & optional & explicit graphs, Laplacians, quick prototypes\\
PyBullet, gym-pybullet-drones \cite{panerati2021learning} & optional & physics simulation of quadrotors for the capstone\\
\bottomrule
\end{tabularx}
\end{table}

The training plan behind this book (\cref{ch:appA}) prescribes a lean stack, listed in \cref{tab:ch02-stack}. The code is one file per chapter, \code{Overleaf/code/chNN\_slug.py} (\code{ch02\_toolbox.py} here, \code{ch04\_astar.py} for \astar), self-contained apart from NumPy and the occasional import from an earlier chapter, and ending in an \code{if \_\_name\_\_ == "\_\_main\_\_":} block that runs a self-test of \code{assert} statements in a few seconds; every number quoted in a worked example comes from running it.

\Cref{lst:ch02-grid-neighbors} shows the successor function of the toolbox grid, the one function that the searches of \cref{ch:ch03,ch:ch04,ch:ch05,ch:ch06} call most often; note the corner-cutting test and the two edge costs of \cref{def:ch02-grid}. In the same file \code{tangent\_points} transcribes \cref{eq:ch02-tangent-points} of \cref{thm:ch02-tangents} line for line, \code{time\_of\_closest\_approach} transcribes \cref{eq:ch02-tca} and \code{time\_to\_collision} is a one-line call to \code{ray\_circle\_intersection} in the relative frame. The class \code{LazyPQ} in \code{code/ch02\_toolbox.py} implements \cref{alg:ch02-lazypq} line for line, and counts its pushes and stale pops.

\begin{lstlisting}[caption={The successor function of the grid (from \code{code/ch02\_toolbox.py}).},label={lst:ch02-grid-neighbors}]
    def neighbors(self, cell):
        """Return [(neighbor, cost)] for the free neighbors of ``cell``."""
        x, y = cell
        moves = MOVES4 if self.connectivity == 4 else MOVES8
        result = []
        for dx, dy in moves:
            nxt = (x + dx, y + dy)
            if not self.is_free(nxt):
                continue
            if dx != 0 and dy != 0:                     # diagonal move
                if not (self.is_free((x + dx, y))
                        and self.is_free((x, y + dy))):
                    continue                            # corner cutting
                result.append((nxt, SQRT2))
            else:
                result.append((nxt, 1.0))
        return result
\end{lstlisting}

% ---------------------------------------------------------------------
\section{Where this fits in the drone system}
\label{sec:ch02-drone}

\begin{dronebox}
Each of the four layers of the hybrid architecture in \cref{ch:ch24} is built from tools of this chapter. The \emph{global planner} searches a grid or lattice (\cref{def:ch02-grid,def:ch02-lattice}) of inflated obstacles (\cref{thm:ch02-inflation}) in space-time (\cref{def:ch02-space-time-state}), compares plans by makespan and sum of costs (\cref{def:ch02-makespan,def:ch02-sum-of-costs}), and runs on the lazy priority queue (\cref{alg:ch02-lazypq}). The \emph{prediction layer} describes an intruder by a Gaussian (\cref{def:appB-gaussian}), read as the covariance ellipse of \cref{def:ch02-covariance-ellipse} and propagated through the double-integrator model (\cref{def:ch02-double-integrator}) with Bayes' rule (\cref{eq:appB-bayes}). The \emph{local safety layer} works in the relative frame of two disks (\cref{thm:ch02-disc-minkowski}), decides with the time to collision (\cref{thm:ch02-closest-approach}) whether an intruder is inside the safety horizon, avoids it with the tangent cone (\cref{thm:ch02-tangents}), and respects the stopping distance (\cref{thm:ch02-stopping-distance}). The \emph{replanning layer} re-searches the grid from the current space-time state with hashed states. In the twelve-week study plan of \cref{ch:appA} this chapter accompanies Week~1, the foundations week, and is the reference to return to in every later week.
\end{dronebox}

% ---------------------------------------------------------------------
\section{Summary}
\begin{summary}
\begin{itemize}
  \item A planner searches a weighted graph through a successor function; grids and lattices are implicit graphs with move costs $1$, $\sqrt2$, $\sqrt3$ and a corner-cutting rule (\cref{def:ch02-graph,def:ch02-grid,def:ch02-lattice}).
  \item A disk or ball robot becomes a point when the obstacles are inflated by its radius: $\Cobs=\mathcal{O}\oplus\disc{\vect{0}}{r}$ (\cref{thm:ch02-inflation}); a drone is such a ball, radius plus margin.
  \item Discrete time adds the wait action and turns the map into the time-expanded graph; a time-indexed path fixes where and when, a plan is one such path per agent, a trajectory is continuous motion (\cref{def:ch02-time-indexed-path,def:ch02-plan,def:ch02-trajectory}).
  \item Path length, travel time, makespan $\makespan=\max_i T_i$, sum of costs $\sumcost=\sum_i T_i$ and minimum separation are the cost measures of the book; a plan without grid conflicts can still bring physical agents too close.
  \item The single integrator chooses velocities, the double integrator chooses accelerations with exact discrete matrices $\mat{A},\mat{B}$; a quadrotor is nearly the latter; the stopping distance is $v^2/(2a_{\max})$.
  \item Distance to a segment, ray--circle intersection, the Minkowski disk, tangent lines and the time of closest approach $t^*=(\pos_{\mathrm{rel}}\cdot\vel_{\mathrm{rel}})/\norm{\vel_{\mathrm{rel}}}^2$ with $d_{\min}=\abs{\pos_{\mathrm{rel}}\times\vel_{\mathrm{rel}}}/\norm{\vel_{\mathrm{rel}}}$ are the geometry of collision avoidance.
  \item A covariance is an ellipse with axes $\sqrt{\lambda_i}\,\vect{e}_i$ (\cref{thm:ch02-ellipse-axes}); in 2D the $1\sigma$ ellipse holds only $39\%$ of the mass; Bayes' rule turns a prior and a measurement into a posterior.
  \item Searches use a binary heap with lazy deletion (push duplicates, skip stale pops) and dictionaries keyed by integer tuples; \python and NumPy are enough for every instance in this book.
\end{itemize}
\end{summary}

\paragraph{Further reading.}
Graphs, shortest paths, heaps and big-O notation are covered in \textcite{cormen2009clrs}, and the search view of planning in \textcite{russell2020aima}. The configuration-space construction and the inflation of \cref{thm:ch02-inflation} go back to \textcite{lozanoperez1983spatial}. Configuration space, Minkowski sums, holonomic and non-holonomic systems and motion models are the subject of \textcite{lavalle2006planning} and \textcite{choset2005principles}; the latter also contains the geometry of distance computations. The definitions of time-indexed paths, makespan and sum of costs follow \textcite{stern2019mapf}. The velocity-obstacle construction that motivates the tangent-line and closest-approach primitives is due to \textcite{fiorini1998motion}. Gaussians, covariance ellipses and Bayes' rule in the form used by tracking are treated in \textcite{thrun2005probabilistic}. For the quadrotor as a double integrator and differential flatness see \textcite{mellinger2011minimum}; for the simulation stack see \textcite{panerati2021learning}.

% ---------------------------------------------------------------------
\section{Exercises}
\label{sec:ch02-exercises}

\begin{exercise}[\difficulty{1}]
\label{exr:ch02-adjacency}
Write the adjacency list of the 8-connected grid graph of the $4\times3$ map \code{["....", ".\#..", "...."]} for the cells $(0,0)$, $(2,1)$ and $(3,0)$, including edge costs. Then show that an empty $W\times H$ grid has $W(H-1)+H(W-1)$ undirected edges when 4-connected and $2(W-1)(H-1)$ more when 8-connected; check the totals $17$ and $29$ for $W=4$, $H=3$ against the self-test of the toolbox.
\end{exercise}

\begin{exercise}[\difficulty{1}]
\label{exr:ch02-diagonal-cost}
Suppose a diagonal move cost $1$ instead of $\sqrt2$. Give a start and goal on an empty 8-connected grid for which the cheapest grid path is then not the shortest polyline through cell centers, and explain what this would do to a heuristic that estimates the straight-line distance (\cref{ch:ch04}).
\end{exercise}

\begin{exercise}[\difficulty{2}]
\label{exr:ch02-inflation}
A single cell of a large grid is blocked. Using the cell-center rule of \code{Grid.inflate}, determine by hand which cells are blocked after inflating by $r=0.6$, $r=1.0$ and $r=1.6$, and count them in each case. At which values of $r$ does the count change next? Compare the blocked-cell count with the area $1+4r+\pi r^2$ of the exact Minkowski sum of the unit square with the disk of radius $r$, and check that the ratio tends to $1$.
\end{exercise}

\begin{exercise}[\difficulty{3}]
\label{exr:ch02-time-expanded}
Let $G=(V,E)$ be undirected. Show that its time-expanded graph with horizon $T_{\max}$ has $\abs{V}(T_{\max}+1)$ vertices and $T_{\max}(\abs{V}+2\abs{E})$ edges, that it is acyclic, and that its paths from $(s,0)$ to $(\gamma,T_{\max})$ correspond one-to-one to the time-indexed paths of \cref{def:ch02-time-indexed-path} with horizon $T_{\max}$.
\end{exercise}

\begin{exercise}[\difficulty{3}]
\label{exr:ch02-costs}
Compute the makespan and the sum of costs of the plan in \cref{ex:ch02-costs}. Then construct an instance with three agents on a corridor in which the plan with the smallest makespan does not have the smallest sum of costs, and vice versa, and prove that no plan is optimal for both objectives.
\end{exercise}

\begin{exercise}[\difficulty{2}]
\label{exr:ch02-stopping}
A drone flies at $v=2$~m/s with $a_{\max}=1$~m/s$^2$ and $\dt=0.1$~s. How many steps of full braking does it need, and how far does it travel according to the exact model \cref{eq:ch02-double-integrator}? Repeat with the forward-Euler update $\pos_{k+1}=\pos_k+\dt\vel_k$, $\vel_{k+1}=\vel_k+\dt\acc_k$ and explain the difference. Finally, write the rule ``a velocity $v$ is admissible if the drone can still stop before the nearest obstacle at distance $d$'' as an inequality; this rule returns in \cref{ch:ch14}.
\end{exercise}

\begin{exercise}[\difficulty{2}]
\label{exr:ch02-closest-approach}
Derive \cref{eq:ch02-tca} by minimizing $D(t)^2$, show that $d_{\min}=\abs{\pos_{\mathrm{rel}}\times\vel_{\mathrm{rel}}}/\norm{\vel_{\mathrm{rel}}}$ when the minimum is not clamped, and derive the time to collision \cref{eq:ch02-ttc}. Evaluate all three for $\pos_{\mathrm{rel}}=(4,2)$, $\vel_{\mathrm{rel}}=(1,1)$ and collision distances $R=1.5$ and $R=1$.
\end{exercise}

\begin{exercise}[\difficulty{2}]
\label{exr:ch02-tangents}
For $\vect{p}=(0,0)$, $\vect{c}=(4,0)$ and $r=2$ compute the tangent points, the tangent length and the half-angle $\theta$ of the cone from \cref{thm:ch02-tangents}. Show in general that $\theta=\arcsin(r/d)$, and describe what happens to the cone as $d\to r^{+}$ and as $d\to\infty$. Why does this matter for a drone that is close to another drone?
\end{exercise}

\begin{exercise}[\difficulty{2}]
\label{exr:ch02-ellipse}
For $\mat{\Sigma}=\begin{psmallmatrix}2&1.2\\1.2&1\end{psmallmatrix}$ compute the eigenvalues, the eigenvectors and the orientation of the major axis by hand, and give the semi-axes of the $1\sigma$ and $2\sigma$ ellipses. Then compute the probability mass inside each ellipse and compare it with the one-dimensional $68\%$ and $95\%$ figures.
\end{exercise}

\begin{exercise}[\difficulty{3} Coding]
\label{exr:ch02-coding}
Extend the toolbox to three dimensions and measure the priority queue. (a)~Write a class \code{Grid3D} with 6- and 26-connectivity, the move costs of \cref{def:ch02-lattice}, the box version of the corner-cutting rule and an \code{inflate} method that samples the Minkowski sum with a ball. (b)~Implement \code{segment\_sphere\_intersects} and the 3D versions of \code{time\_of\_closest\_approach} and \code{time\_to\_collision}, and test them with \code{assert}s against hand-computed values, as the self-test of \code{ch02\_toolbox.py} does. (c)~Using \code{space\_time\_successors} and \code{LazyPQ}, write a uniform-cost search over space-time states on a 2D grid (\cref{ch:ch03} explains why it is correct), run it on a $50\times50$ map with random obstacles, and report the pushes and stale pops; explain why the stale pops never exceed the pushes.
\end{exercise}
